% !TeX root=../main.tex
\chapter{بحث و نتیجه‌گیری}
%\thispagestyle{empty} 

% sec 1
\section{مقدمه}
\label{sec:chapter5_intro}

این فصل به جمع‌بندی نهایی پژوهش حاضر و تبیین برداشت‌های حاصل از یافته‌های آن اختصاص دارد.
در فصل‌های پیشین، مسئلهٔ پژوهش، پیشینهٔ علمی مرتبط، روش پیشنهادی و نتایج تجربی
به‌صورت گام‌به‌گام ارائه و تحلیل شدند.
در این مرحله، هدف اصلی پاسخ‌گویی به این پرسش بنیادین است که
یافته‌های به‌دست‌آمده چه معنایی دارند
و چه ارزش علمی و عملی برای حوزهٔ تخمین عمق تک‌دوربینه\enfootnote{Monocular Depth Estimation} ایجاد می‌کنند.

در فصل نخست، ضرورت پرداختن به مسئلهٔ تخمین عمق متریک\enfootnote{Metric Depth}
و چالش‌های ذاتی آن، به‌ویژه ابهام مقیاس\enfootnote{Scale Ambiguity} و وابستگی به پارامترهای دوربین، مطرح شد.
در ادامه، مرور ادبیات (فصل دوم) نشان داد که
جریان غالب روش‌های موجود،
یا به عمق غیرمتریک\enfootnote{Non-metric Depth} بسنده کرده‌اند
و یا نقش هندسهٔ تصویربرداری را به‌صورت ضمنی در نظر گرفته‌اند~\cite{Eigen2014Depth,Godard2019Digging,Ranftl2019Robust,DepthAnythingV2}؛
در برابر آن، خط پژوهشی موازی‌ای وجود دارد که پارامترهای دوربین را صریحاً
وارد مدل می‌کند~\cite{Facil2019CAMConvs,Yin2023Metric3D,Saxena2023FOV,Piccinelli2024UniDepth}،
اما در هیچ‌یک از آن‌ها ورود $K$ از تغییر هم‌زمان معماری، داده و برنامهٔ آموزش
تفکیک نشده است. همین تفکیک، هدف پژوهش حاضر قرار گرفت.
در فصل سوم، چارچوب روش‌شناختی پژوهش ارائه شد
و نشان داده شد که چگونه می‌توان
با ادغام صریح ماتریس مشخصهٔ دوربین\enfootnote{Camera Intrinsic Matrix}
در معماری یک مدل یادگیری عمیق،
رابطهٔ میان تصویر، هندسهٔ دوربین و عمق متریک را
به‌صورت مستقیم مدل‌سازی کرد.
در فصل چهارم نیز،
نتایج کمی و کیفی حاصل از آزمایش‌ها
برای ارزیابی صحت این رویکرد گزارش و تحلیل شدند.

اکنون در فصل حاضر،
با تکیه بر نتایج عملی استخراج‌شده
و بدون ارائهٔ داده‌ها یا آزمایش‌های جدید،
برداشت نهایی از یافته‌های پژوهش مطرح می‌شود.
در این فصل، تلاش می‌شود
ضمن جمع‌بندی دستاوردهای اصلی،
میزان تحقق اهداف پژوهش مشخص شود،
نوآوری‌های تئوریک و عملی کار تبیین گردد،
و محدودیت‌ها و مسیرهای آتی پژوهش
در چارچوب مسئلهٔ مورد مطالعه بیان شوند.
بدین ترتیب، این فصل نقش حلقهٔ پایانی تز را ایفا می‌کند
و ارتباط منطقی میان مسئلهٔ اولیه،
روش انجام تحقیق
و نتایج به‌دست‌آمده را به‌صورت یکپارچه روشن می‌سازد.


% sec 2
\section{جمع‌بندی یافته‌های اصلی پژوهش}
\label{sec:chapter5_summary}

هدف اصلی این پژوهش،
بررسی امکان دستیابی به تخمین عمق متریک پایدار
با استفاده از یک مدل بنیادین تخمین عمق تک‌دوربینه
و تحلیل نقش اطلاعات هندسی دوربین
در تحقق این هدف بود.
یافته‌های به‌دست‌آمده در فصل‌های سوم و چهارم
به‌صورت منسجم نشان می‌دهند که
ابهام مقیاس، نه یک ضعف پیاده‌سازی،
بلکه یک محدودیت هندسی در مدل‌هایی است
که صرفاً بر اطلاعات تصویری تکیه دارند --- محدودیتی که در رژیم دادهٔ محدود،
عملاً غیرقابل جبران است. لازم است تصریح شود که این محدودیت، \emph{مطلق}
نیست: مدل‌هایی نظیر \lr{Depth Pro} نشان داده‌اند که با مقیاس دادهٔ به‌مراتب
بزرگ‌تر می‌توان فاصلهٔ کانونی را از خود تصویر تخمین زد و بدون فراداده به عمق
متریک رسید~\cite{Bochkovskii2025DepthPro}. آنچه یافته‌های این پژوهش نشان
می‌دهد، ناممکن‌بودن بازیابی مقیاس بدون $K$ نیست، بلکه پرهزینه‌بودن آن است.

نتایج نشان داد که
مدل‌های متداول تخمین عمق تک‌دوربینه،
از جمله \basemodel{} مورد بررسی در این پژوهش،
توانایی یادگیری ساختار نسبی عمق صحنه را دارند،
اما در غیاب اطلاعات صریح مربوط به هندسهٔ دوربین،
قادر به بازیابی مقیاس فیزیکی پایدار نیستند.
این مسئله به‌صورت ناپایداری خروجی عمق
در برابر تغییر دیتاست و پارامترهای تصویربرداری
در معیارهای متریک آشکار شد،
در حالی که معیارهای نرمال‌شده
در بسیاری از موارد این ناپایداری را پنهان می‌کنند.

در مقابل،
یافته‌های این پژوهش نشان می‌دهد که
ادغام صریح ماتریس مشخصهٔ دوربین
در فرآیند تخمین عمق
می‌تواند این محدودیت ذاتی را تا حد قابل توجهی کاهش دهد.
\proposedmodel{} با بهره‌گیری از اطلاعات هندسی دوربین،
وابستگی خروجی عمق به شرایط تصویربرداری را کاهش داده
و به تخمینی پایدارتر از نظر مقیاس فیزیکی دست یافته است.
این بهبود نه‌تنها در ارزیابی‌های کمی،
بلکه در تحلیل‌های کیفی نیز
به‌صورت یکنواختی بهتر سطوح،
حفظ مقیاس اجسام
و بازنمایی واقع‌بینانه‌تر نواحی دوردست مشاهده شد.

مطالعهٔ کنترل‌شدهٔ نهایی بر روی پروتکل استاندارد \lr{NYUv2}
(بخش~\ref{sec:controlled_nyu_study})، این یافته‌ها را با دقت آماری مشخص‌تری صورت‌بندی کرد.
\proposedmodel{} در پیکربندی \lr{ViT-L}،
\publishedbasemodel{} را
در تمامی معیارهای ارزیابی و با معناداری آماری بالا پشت سر گذاشت:
در پروتکل نسبی، کاهش AbsRel از \fadecimal{۰٫۰۴۲۵} به \fadecimal{۰٫۰۳۷۷}
($p=2.5\times10^{-8}$، آزمون t زوجی، $n=654$)
و در ارزیابی متریک بدون تراز،
کاهش AbsRel از \fadecimal{۰٫۲۰۷} به \fadecimal{۰٫۰۷۹} (بهبود \fadecimal{۲٫۶} برابری).
در عین حال، مقایسه با \controlmodel{} هم‌داده نشان داد که
سهم خالص ورودی پارامترهای دوربین،
وابسته به پروتکل ارزیابی است:
در ارزیابی نسبی که تراز مقیاسِ به‌ازای هر تصویر،
اطلاعات مقیاس را حذف می‌کند،
این سهم تنها در انکودر کوچک معنادار بود؛
اما در ارزیابی متریک،
برتری \proposedmodel{} در دو اجرای مستقل تکرار شد
(RMSE با $p=4.9\times10^{-3}$ و $p=6.5\times10^{-5}$)
و با افزایش تنوع میدان دید در داده‌های آموزشی،
به‌صورت نظام‌مند بزرگ‌تر گردید.
نقطهٔ اوج این الگوی پاسخ به دوز،
آموزش با افزونه‌سازی لرزش کانونی بود
(زیربخش~\ref{subsec:focal_jitter_results})
که مقیاس مطلق را بدون $K$ اساساً غیرقابل بازیابی می‌کند:
در این پیکربندی، \proposedmodel{}
در تمامی معیارها با معناداری در مرتبهٔ $p \le 10^{-24}$
از \controlmodel{} هم‌داده و هم‌بذر پیشی گرفت
(AbsRel: \fadecimal{۰٫۰۸۵} در برابر \fadecimal{۰٫۱۰۳})
و آزمون سازوکاری جاروی $K$ نشان داد که
خروجی مدل تقریباً دقیقاً مطابق معادلهٔ روزنه‌ای
با فاصلهٔ کانونی مقیاس می‌شود
(شیب لگاریتمی \fadecimal{۰٫۹۳} در برابر مقدار ایده‌آل ۱؛
\controlmodel{}: صفر).
مطالعهٔ فروکاست، عامل علّی را به‌طور قطعی
به خودِ افزونه‌سازی لرزش کانونی منتسب کرد
(بازتولید کامل اثر با $\lambda$ بدون تغییر:
AbsRel \fadecimal{۰٫۰۸۳} در برابر \fadecimal{۰٫۱۰۷} برای \controlmodel{}، $p \le 10^{-39}$،
شیب جاروی \fadecimal{۰٫۹۷})
و آزمون حذف و تحریف $K$ نشان داد که
این اتکا علّی است:
حذف $K$ دقت را ۵۹ درصد تخریب می‌کند
و تحریف عمدی فاصلهٔ کانونی،
خطا را دقیقاً به اندازه‌ای که هندسهٔ روزنه‌ای
پیش‌بینی می‌کند افزایش می‌دهد.
مجموعهٔ این شواهد --- الگوی پاسخ به دوز در چهار زوج آموزشی مستقل،
برتری همه‌جانبهٔ آماری،
و انطباق کمی رفتار مدل با هندسهٔ تصویربرداری ---
سودمندی علّی اطلاعات هندسی دوربین را
به‌صورت هم‌گرا اثبات می‌کند.

همچنین،
تحلیل بین‌دیتاستی نتایج نشان داد که
پایداری متریک \proposedmodel{}
در محیط‌های متنوع داخلی و خارجی حفظ می‌شود،
هرچند کیفیت و ماهیت داده‌های عمق مرجع
می‌تواند بر مقادیر مطلق خطا اثرگذار باشد.
مطالعهٔ موردی دیتاست CityScapes
نشان داد که تفسیر نتایج
باید با در نظر گرفتن محدودیت‌های ذاتی هر دیتاست انجام شود
و مقایسهٔ مستقیم خطاها
بدون توجه به کیفیت داده‌های مرجع
می‌تواند گمراه‌کننده باشد.

در مجموع،
یافته‌های این پژوهش نشان می‌دهد که
دستیابی به عمق متریک پایدار
در تخمین عمق تک‌دوربینه
نه از طریق پیچیده‌تر کردن معماری شبکه،
بلکه از طریق هم‌راستا کردن مدل یادگیری
با اصول هندسهٔ تصویربرداری امکان‌پذیر است.
این نتیجه،
چارچوب نظری و عملی ارائه‌شده در این پژوهش را تأیید می‌کند
و مبنایی برای تفسیر نوآوری‌ها،
محدودیت‌ها و مسیرهای آتی تحقیق
در بخش‌های بعدی این فصل فراهم می‌سازد.

\begin{figure}[htbp]
    \centering
    \begin{tikzpicture}[font=\small]
    \tikzset{
      stg/.style={thblock, minimum height=1.7cm, text width=3.15cm, font=\scriptsize},
      ev/.style={thdata, minimum height=1.15cm, text width=3.15cm, font=\scriptsize},
    }
    %%% زنجیرهٔ اصلی (راست به چپ)
    \node[stg, fill=thPanelRed, draw=thControl] (gap) at (10.8,1.7)
        {\rl{\textbf{چالش شناسایی‌شده}}\\[2pt] \rl{ابهام مقیاس یک محدودیت \emph{ذاتی هندسی} است، نه ضعف معماری}};
    \node[stg, fill=thPanelAlt, draw=thAccent] (method) at (7.2,1.7)
        {\rl{\textbf{راهکار پیشنهادی}}\\[2pt] \rl{تزریق دروازه‌بانی‌شدهٔ ماتریس $K$ به‌صورت یک توکن به انکودر}};
    \node[stg, fill=thPanel, draw=thProposed] (design) at (3.6,1.7)
        {\rl{\textbf{طراحی آزمایش}}\\[2pt] \rl{زوج هم‌داده و هم‌بذر، ارزیابی متریک بدون تراز، لرزش کانونی}};
    \node[stg, fill=thPanelGreen, draw=thIdeal] (result) at (0.0,1.7)
        {\rl{\textbf{دستاورد}}\\[2pt] \rl{عمق متریک پایدار با اتکای علّی اثبات‌شده به $K$}};

    \draw[tharrow] (gap) -- (method);
    \draw[tharrow] (method) -- (design);
    \draw[tharrow] (design) -- (result);

    %%% شواهد پشتیبان
    \node[ev, fill=white] (e1) at (7.2,-0.95)
        {\rl{الگوی پاسخ به دوز در چهار زوج مستقل}};
    \node[ev, fill=white] (e2) at (3.6,-0.95)
        {\rl{شیب جاروی $K$: از $0.06$ به $0.97$}};
    \node[ev, fill=white] (e3) at (0.0,-0.95)
        {\rl{افت ۵۹ درصدی دقت با حذف $K$}};

    \draw[thdashed] (design.south) -- (e1.north);
    \draw[thdashed] (design.south) -- (e2.north);
    \draw[thdashed] (result.south) -- (e3.north);
    \node[thnote, anchor=east] at (10.9,-0.95) {\rl{شواهد پشتیبان:}};
    \end{tikzpicture}
    \caption[نمای کلی از فرآیند و دستاوردهای پژوهش]{نمای کلی از فرآیند
    پژوهش: حرکت از شناسایی یک محدودیت ذاتی هندسی در مدل‌های موجود، به یک
    راهکار هندسه‌آگاه، سپس به طرح آزمایشی که سنجش منصفانهٔ آن راهکار را ممکن
    می‌سازد، و در نهایت به دستاورد پایداری متریک. سه شاهد پشتیبان در ردیف
    پایین، همان یافته‌هایی هستند که ادعای این زنجیره را از سطح هم‌بستگی به
    سطح علیت ارتقا می‌دهند.}
    \label{fig:contributions_summary}
\end{figure}



% sec 3
\section{تفسیر نهایی نتایج در چارچوب هندسهٔ تصویربرداری}
\label{sec:chapter5_geometric_interpretation}

یافته‌های این پژوهش را می‌توان به‌صورت منسجم
در چارچوب هندسهٔ تصویربرداری و مدل فیزیکی دوربین تفسیر کرد.
در مدل سوراخ‌سوزنی، رابطهٔ میان مختصات تصویر و عمق صحنه
به‌صورت مستقیم به پارامترهای کانونی و هندسهٔ پروجکشن وابسته
است~\cite{HartleyZisserman2004}.
از این رو، عمق متریک ذاتاً یک کمیت هندسی است
و بازیابی آن \emph{تنها از روی هندسهٔ یک تصویر} و بدون اطلاع از پارامترهای
دوربین، به‌صورت یکتا امکان‌پذیر نیست. آنچه مدل‌های داده‌محور در عمل انجام
می‌دهند، جایگزین‌کردن این اطلاعات غایب با آمار سرنخ‌های تصویری (اندازهٔ ظاهری
اجسام آشنا، هندسهٔ خطوط گریز و مانند آن) است؛ راهبردی که با مقیاس دادهٔ کافی
تا حد قابل توجهی کارآمد است~\cite{Bochkovskii2025DepthPro} اما در رژیم
دادهٔ محدودِ این پژوهش، ناپایدار می‌ماند.

نتایج تجربی این پژوهش نشان دادند که
مدل‌هایی که تنها بر اطلاعات تصویری تکیه دارند،
در بهترین حالت قادر به یادگیری ساختار نسبی عمق هستند.
این رفتار با تحلیل‌های نظری فصل سوم هم‌خوانی دارد،
زیرا در غیاب ماتریس مشخصهٔ دوربین،
نگاشت تصویر به عمق دارای درجه‌ای از آزادی مقیاس است
که شبکهٔ یادگیری عمیق نمی‌تواند آن را
صرفاً از دادهٔ تصویری حذف کند.
به همین دلیل،
پایداری مشاهده‌شده در معیارهای غیرمتریک
لزوماً به معنای بازیابی صحیح عمق فیزیکی نیست.

در مقابل،
زمانی که اطلاعات هندسی دوربین
به‌صورت صریح در فرآیند تخمین عمق لحاظ می‌شود،
مدل یادگیری با محدودیت‌های فیزیکی سیستم تصویربرداری
هم‌راستا می‌گردد.
نتایج این پژوهش نشان می‌دهد که
ادغام صریح ماتریس مشخصهٔ دوربین
در عمل نقش یک قید هندسی را ایفا می‌کند
که فضای پاسخ‌های ممکن شبکه را محدود می‌سازد.
این محدودسازی باعث می‌شود که
مدل به‌جای اتکا به هم‌بستگی‌های آماری ناپایدار،
وابستگی‌های هندسی معنادار میان تصویر و عمق را بیاموزد.

تحلیل رفتار معیارهای ارزیابی نیز
از این تفسیر هندسی پشتیبانی می‌کند.
معیارهای نرمال‌شده، به‌دلیل حذف ضمنی اثر مقیاس،
نمی‌توانند تغییرات ناشی از پارامترهای دوربین را به‌خوبی منعکس کنند،
در حالی که معیارهای متریک
مستقیماً به روابط پروجکشن وابسته‌اند
و نسبت به خطاهای هندسی حساس هستند.
بهبود پایداری در این معیارها
نشان‌دهندهٔ آن است که
\proposedmodel{} توانسته است
وابستگی خروجی عمق به هندسهٔ تصویربرداری را
به‌صورت سازگارتر مدل‌سازی کند.

از منظر هندسی،
یافته‌های این پژوهش نشان می‌دهند که
در رژیم فاین‌تیون با دادهٔ محدود، یادگیری عمیق بدون لحاظ صریح مدل
تصویربرداری، جایگزین کارآمدی برای قیود فیزیکی نیست؛ همان نتیجه‌ای که کارهای
شرطی‌شده بر هندسهٔ دوربین نیز --- از راه‌های متفاوت --- به آن
رسیده‌اند~\cite{Facil2019CAMConvs,Yin2023Metric3D,Piccinelli2024UniDepth}.
در عوض،
ترکیب آگاهانهٔ یادگیری داده‌محور
با مدل‌های کلاسیک هندسی
می‌تواند به تخمین‌هایی منجر شود
که هم از نظر بصری معنادارند
و هم از نظر فیزیکی پایدار.
این تفسیر نشان می‌دهد که
دستاورد اصلی پژوهش حاضر
نه صرفاً بهبود یک مدل خاص،
بلکه ارائهٔ شواهد تجربی
برای اهمیت ادغام هندسهٔ تصویربرداری
در طراحی مدل‌های تخمین عمق تک‌دوربینه است.


% sec 4
\section{نوآوری‌های تحقیق}
\label{sec:chapter5_novelty}

بر اساس یافته‌های نظری و تجربی ارائه‌شده در این پژوهش،
نوآوری‌های تحقیق حاضر را می‌توان
در دو سطح تئوریک و عملی دسته‌بندی کرد.
این نوآوری‌ها نه بر مبنای ادعاهای کلی،
بلکه بر اساس نتایج عملی استخراج‌شده
و تحلیل‌های انجام‌شده در فصل‌های سوم و چهارم
تشریح می‌شوند.

\subsection{نوآوری تئوریک}
\label{subsec:theoretical_novelty}

نوآوری تئوریک اصلی این پژوهش
در بازتعریف مسئلهٔ تخمین عمق تک‌دوربینه
در چارچوب هندسهٔ تصویربرداری نهفته است.
در جریان غالب کارهای پیشین،
تخمین عمق به‌عنوان یک مسئلهٔ صرفاً داده‌محور
مدل‌سازی شده
و نقش پارامترهای هندسی دوربین
یا نادیده گرفته شده
یا به‌صورت ضمنی در داده‌ها فرض شده است~\cite{Eigen2014Depth,Godard2019Digging,Ranftl2021DPT,DepthAnythingV2}.

صورت‌بندی هندسی ابهام مقیاس، به‌خودی‌خود تازه نیست و پیش‌تر در
\lr{Metric3D} به‌عنوان انگیزهٔ «فضای دوربین متعارف» ارائه شده
است~\cite{Yin2023Metric3D}. سهم نظری پژوهش حاضر، \emph{اندازه‌گیری} این
ابهام است: با ثابت نگه‌داشتن همهٔ عوامل جز دسترسی به $K$، نشان داده شد که
بخشی از خطای متریک که پیش‌تر به کمبود ظرفیت یا کمبود داده نسبت داده می‌شد،
در واقع سهم اطلاعات هندسیِ غایب است، و این سهم با آزمون‌های جاروی $K$ و حذف
و تحریف $K$، به‌صورت کمی قابل ردیابی است.

بر این اساس، ادعای قابل دفاع این پژوهش چنین است: در رژیم فاین‌تیون با دادهٔ
محدود و ستون فقرات فریزشده، بازیابی مقیاس متریک پایدار \emph{بدون} ورودی
صریح هندسهٔ دوربین به‌شدت پرهزینه و ناپایدار است. این ادعا با ادعای قوی‌تر
«ناممکن‌بودن» یکسان نیست؛ کارهایی با مقیاس دادهٔ بسیار بزرگ‌تر نشان داده‌اند
که فاصلهٔ کانونی را می‌توان از خود تصویر بازسازی کرد~\cite{Bochkovskii2025DepthPro}.

افزون بر این،
پژوهش حاضر با تفکیک روشن
میان معیارهای متریک و غیرمتریک
و تحلیل رفتار متفاوت آن‌ها
نشان می‌دهد که
ارزیابی‌های متداولِ مبتنی بر معیارهای نرمال‌شده --- که از پروتکل
\lr{Eigen}~\cite{Eigen2014Depth} و تراز مقیاس و جابجایی خانوادهٔ
\lr{MiDaS}~\cite{Ranftl2019Robust} ریشه می‌گیرند و در ارزیابی مدل‌های
بنیادین اخیر نیز رایج‌اند~\cite{DepthAnythingV2} ---
می‌توانند تصویری ناقص یا حتی گمراه‌کننده
از عملکرد واقعی مدل‌ها ارائه دهند.
این دیدگاه،
چارچوبی مفهومی برای بازنگری
در نحوهٔ ارزیابی مدل‌های تخمین عمق
در ادبیات موجود فراهم می‌کند.

\subsection{نوآوری عملی}
\label{subsec:practical_novelty}

در سطح عملی،
نوآوری این پژوهش
در ارائهٔ یک راهکار ساده، قابل پیاده‌سازی
و سازگار با معماری‌های موجود
برای بهبود تخمین عمق متریک است.
برخلاف روش‌هایی که ورود اطلاعات دوربین را با تغییر گستردهٔ معماری یا
بازآموزی در مقیاس بزرگ همراه می‌کنند --- از کانال‌های مختصاتی در
\lr{CAM-Convs}~\cite{Facil2019CAMConvs} تا شاخهٔ دوربین درون‌مدلی
\lr{UniDepth}~\cite{Piccinelli2024UniDepth} ---،
روش پیشنهادی با حفظ هستهٔ معماری پایه
و افزودن صریح اطلاعات ماتریس مشخصهٔ دوربین
در قالب یک توکن دروازه‌بانی‌شده به انکودر شبکه،
توانسته است پایداری متریک خروجی عمق را افزایش دهد.

یافته‌های عملی این پژوهش
نشان می‌دهد که
حتی بدون تغییر انکودر
و با ثابت نگه‌داشتن وزن‌های استخراج ویژگی،
می‌توان با بهره‌گیری از اطلاعات هندسی مناسب،
به بهبود معنادار در معیارهای متریک دست یافت.
این نتیجه،
توصیه‌ای عملی برای پژوهشگران و کاربردهای صنعتی است
که نشان می‌دهد
نادیده گرفتن اطلاعات دوربین
می‌تواند منجر به نتایج ناپایدار
در کاربردهایی شود
که به مقیاس فیزیکی دقیق نیاز دارند.

در مجموع،
نوآوری عملی این پژوهش
در ارائهٔ شواهد تجربی است
که نشان می‌دهد
ادغام آگاهانهٔ هندسهٔ دوربین
در مدل‌های یادگیری عمیق
می‌تواند بدون افزایش قابل توجه پیچیدگی سیستم،
کیفیت و پایداری تخمین عمق متریک را بهبود بخشد.

\subsection{دستاوردهای روش‌شناختی و تشخیصی}
\label{subsec:methodological_contributions}

در کنار نوآوری‌های فوق، این پژوهش مجموعه‌ای از دستاوردهای روش‌شناختی
دارد که اعتبار نتایج را تقویت کرده و برای پژوهش‌های مشابه قابل استفاده‌اند:

\begin{itemize}
    \item \textbf{شناسایی و رفع نشت داده در تقسیم‌بندی \lr{NYUv2}:}
    نشان داده شد که تقسیم تصادفی متداول بر روی مجموعهٔ برچسب‌دار،
    بیش از ۸۰ درصد تصاویر تقسیم آزمون استاندارد \lr{Eigen}~\cite{Eigen2014Depth}
    را وارد آموزش می‌کند؛
    اصلاح این نشت، پیش‌شرط هر مقایسهٔ معتبر با اعدادی است که مدل‌های مرجع
    بر همین تقسیم گزارش می‌کنند~\cite{DepthAnythingV2}.

    \item \textbf{تشخیص و تبیین یک شکست خاموش در فاین‌تیون مدل‌های نسبی:}
    ناهم‌خوانی فضای هدف آموزش (عمق مستقیم) با فضای خروجی مدل
    ازپیش‌آموزش‌دیده (عمق معکوس)، در تعامل با جزئیات تابع هزینهٔ
    تغییرناپذیر نسبت به مقیاس، به فروریزش برگشت‌ناپذیر سر شبکه
    می‌انجامد در حالی که منحنی هزینه طبیعی می‌نماید.
    ریشه‌یابی و رفع این پدیده، آموزش موفق مدل‌های بزرگ را ممکن ساخت.

    \item \textbf{افزونه‌سازی هندسی آگاه از دوربین:}
    راهبردی که با برش‌های مقیاس‌متغیر و به‌روزرسانی دقیق ماتریس $K$،
    از یک مجموعه‌دادهٔ تک‌دوربینه، داده‌هایی با میدان دید متنوع و
    برچسب هندسی صحیح می‌سازد. این راهبرد به‌تنهایی
    بهترین نتایج پروتکل نسبی این پژوهش را رقم زد
    و بستر آزمون علّی سهم پارامترهای دوربین را فراهم کرد.

    \item \textbf{تحلیل وابستگی نتیجه‌گیری به پروتکل ارزیابی:}
    نشان داده شد که پروتکل ارزیابی نسبی، به‌صورت ساختاری
    نسبت به سهم اصلی پارامترهای دوربین (مقیاس مطلق) نابیناست
    و سنجش منصفانهٔ این سهم، مستلزم ارزیابی متریک بدون تراز است.
    این یافته برای طراحی ارزیابی در پژوهش‌های آتی حوزه اهمیت دارد.

    \item \textbf{افزونه‌سازی لرزش کانونی:}
    تبدیل دقیق $(f, Z) \mapsto (\alpha f, \alpha Z)$ برگرفته از معادلهٔ
    روزنه‌ای که بدون هیچ بازنمونه‌برداری، بازیابی مقیاس را برای مدل
    فاقد $K$ به مسئله‌ای بدساخت تبدیل می‌کند و اتکای مدل به پارامترهای
    دوربین را از یک گزینه به یک ضرورت آموزشی ارتقا می‌دهد
    (زیربخش~\ref{subsec:focal_jitter}).

    \item \textbf{آزمون تشخیصی جاروی $K$:}
    سنجه‌ای سازوکاری (شیب لگاریتمی پاسخ عمق به مقیاس‌دهی مصنوعی
    فاصلهٔ کانونی) که مستقیماً اندازه می‌گیرد مدل چقدر $K$ را
    «می‌خواند»، و نشان داد که معیارهای تجمیعی و حتی مقدار پارامتر
    دروازهٔ کدگذار دوربین می‌توانند دربارهٔ میزان واقعی مصرف $K$
    گمراه‌کننده باشند (زیربخش~\ref{subsec:intrinsics_variation}).

    \item \textbf{زوج‌سازی کامل آزمایش‌ها با قفل بذر تصادفی:}
    یکسان‌سازی ترتیب داده، وارونه‌سازی‌ها و نمونه‌های افزونه‌سازی
    میان بازوهای \proposedmodel{} و \controlmodel{}، به‌گونه‌ای که تفاوت نتایج تنها به
    متغیر مورد مطالعه (ورودی $K$) قابل انتساب باشد؛ صحت آن با برابری
    بیت‌به‌بیت هزینهٔ گام صفر دو بازو راستی‌آزمایی شد.
\end{itemize}



جدول~\ref{tab:contributions_summary} این نوآوری‌ها را در سه سطح تئوریک،
عملی و روش‌شناختی جمع‌بندی می‌کند و برای هر مورد، شاهد پشتیبان و
قابلیت استفادهٔ مجدد آن در پژوهش‌های دیگر را مشخص می‌سازد.

\begin{table}[htbp]
\centering
\caption[جمع‌بندی نوآوری‌ها و دستاوردهای پژوهش]{جمع‌بندی نوآوری‌ها و
دستاوردهای پژوهش در سه سطح تئوریک، عملی و روش‌شناختی، به‌همراه شاهد پشتیبان
هر مورد.}
\label{tab:contributions_summary}
\footnotesize
\renewcommand{\arraystretch}{1.3}
\begin{tabularx}{\textwidth}{@{}
  >{\raggedleft\arraybackslash}p{0.105\textwidth}
  >{\raggedleft\arraybackslash}p{0.19\textwidth}
  >{\raggedleft\arraybackslash}X
  >{\raggedleft\arraybackslash}p{0.215\textwidth}@{}}
\toprule
سطح & دستاورد & شرح & شاهد پشتیبان \\ \midrule
تئوریک & بازتعریف ابهام مقیاس به‌عنوان محدودیت هندسی &
ابهام مقیاس یک ویژگی ذاتی نگاشت تصویربرداری است، نه ضعف معماری یا کمبود داده؛ بنابراین با بزرگ‌ترکردن شبکه رفع نمی‌شود &
شکل~\ref{fig:scale_ambiguity} و آزمون حذف $K$ (\ref{subsec:k_knockout}) \\
تئوریک & نابینایی ساختاری پروتکل نسبی &
تراز مقیاس به‌ازای هر تصویر، دقیقاً همان کمیتی را حذف می‌کند که $K$ فراهم می‌آورد؛ پس ارزیابی نسبی نمی‌تواند سهم $K$ را بسنجد &
بخش~\ref{subsec:aug_basic_results} \\
عملی & شرطی‌سازی توکنی با دروازهٔ صفرمقدار &
افزودن $K$ بدون تغییر هستهٔ معماری و بدون آسیب به وزن‌های ازپیش‌آموخته، حتی با انکودر کاملاً فریزشده &
شکل~\ref{fig:proposed_architecture} \\
روش‌شناختی & افزونه‌سازی هندسی آگاه از دوربین &
ساخت داده‌های با میدان دید متنوع و برچسب هندسی صحیح از یک مجموعه‌دادهٔ تک‌دوربینه &
شکل~\ref{fig:zoom_crop_aug}؛ بهترین نتایج پروتکل نسبی \\
روش‌شناختی & افزونه‌سازی لرزش کانونی &
تبدیل دقیق $(f,Z)\mapsto(\alpha f,\alpha Z)$ که بازیابی مقیاس را برای مدل فاقد $K$ بدساخت می‌کند، بدون هیچ بازنمونه‌برداری یا نشت اطلاعاتی &
شکل~\ref{fig:focal_jitter} و جدول~\ref{tab:dose_response_summary} \\
روش‌شناختی & آزمون تشخیصی جاروی $K$ &
سنجه‌ای مستقیم برای اینکه مدل چقدر واقعاً $K$ را «می‌خواند»؛ نشان داد معیارهای تجمیعی و حتی مقدار دروازهٔ کدگذار می‌توانند گمراه‌کننده باشند &
شکل~\ref{fig:intrinsics_sensitivity} \\
روش‌شناختی & زوج‌سازی کامل با قفل بذر تصادفی &
یکسان‌سازی ترتیب داده، وارونه‌سازی‌ها و نمونه‌های افزونه‌سازی میان دو بازو، تا تفاوت نتایج تنها به ورودی $K$ قابل انتساب باشد &
برابری بیت‌به‌بیت هزینهٔ گام صفر دو بازو \\
روش‌شناختی & مستندسازی سه شکست خاموش &
نشت تقسیم‌بندی، ناهم‌خوانی فضای هدف، و ناهم‌خوانی معماری سر شبکه؛ هر سه بدون هیچ پیام خطایی رخ می‌دادند &
جدول~\ref{tab:silent_bugs} \\
\bottomrule
\end{tabularx}
\end{table}

% sec 5
\section{محدودیت‌ها و چالش‌های پژوهش}
\label{sec:chapter5_limitations}

با وجود نتایج مثبت و معنادار حاصل از این پژوهش،
لازم است محدودیت‌ها و چالش‌هایی که
در مسیر انجام تحقیق وجود داشته‌اند
به‌صورت شفاف بیان شوند.
ذکر این موارد نه‌تنها به درک دقیق‌تر
دامنهٔ اعتبار یافته‌ها کمک می‌کند،
بلکه مسیر توسعه و بهبود پژوهش‌های آتی
در این حوزه را نیز روشن‌تر می‌سازد.

نخستین محدودیت این پژوهش
به وابستگی روش پیشنهادی به صحت
ماتریس مشخصهٔ دوربین بازمی‌گردد.
در این تحقیق فرض شده است که
پارامترهای درونی دوربین
به‌صورت دقیق و بدون خطا در دسترس هستند.
آزمون حذف و تحریف $K$ (زیربخش~\ref{subsec:k_knockout})
ابعاد این وابستگی را کمی کرده است:
حذف کامل $K$ حدود ۵۹ درصد و تحریف ۳۰ تا ۴۰ درصدیِ
فاصلهٔ کانونی، چند برابر خطای پایه به مدل تحمیل می‌کند.
این رفتار، روی دیگر سکهٔ اتکای علّی مدل به $K$ است:
در عمل، به‌ویژه در کاربردهای واقعی که پارامترها ممکن است
ناقص، تقریبی یا همراه با نویز باشند،
حساسیت مدل به خطای کالیبراسیون باید در طراحی سامانه لحاظ شود؛
مطالعهٔ نظام‌مند رفتار مدل تحت عدم قطعیت‌های واقع‌گرایانهٔ کالیبراسیون
(نه تحریف‌های عمدی بزرگ)
موضوعی برای پژوهش‌های آتی است.

محدودیت دیگر به کیفیت و ماهیت
داده‌های عمق مرجع در دیتاست‌های مختلف مربوط است.
همان‌گونه که در فصل چهارم نشان داده شد،
دیتاست‌ها از نظر نحوهٔ تولید عمق مرجع،
دامنهٔ عمق و یکنواختی داده‌ها
تفاوت‌های اساسی دارند.
این ناهمگنی می‌تواند
بر مقادیر مطلق خطاهای متریک اثرگذار باشد
و تفسیر نتایج را تحت تأثیر قرار دهد.
هرچند در این پژوهش تلاش شده است
این مسئله از طریق تحلیل‌های آماری
و بررسی‌های موردی کنترل شود،
اما محدودیت‌های ذاتی داده‌های مرجع
به‌طور کامل قابل حذف نیستند.

از منظر معماری،
روش پیشنهادی بر پایهٔ یک مدل خاص
از خانوادهٔ \lr{Depth Anything V2}
پیاده‌سازی و ارزیابی شده است.
اگرچه چارچوب ارائه‌شده
به‌صورت مفهومی مستقل از معماری خاص است،
اما تعمیم تجربی نتایج به سایر معماری‌ها
یا انکودرهای متفاوت
در این پژوهش به‌صورت مستقیم بررسی نشده است.
از این رو،
نتایج ارائه‌شده
بیشتر در چارچوب معماری مورد استفاده
قابل تفسیر هستند.

در نهایت،
این پژوهش تمرکز خود را
بر تخمین عمق تک‌دوربینهٔ ایستا
و تصاویر منفرد قرار داده است.
اطلاعات زمانی یا چندنمایی
که می‌توانند قیود هندسی مکملی فراهم کنند،
در چارچوب این تحقیق مورد استفاده قرار نگرفته‌اند.
این انتخاب آگاهانه،
برای تمرکز بر نقش ماتریس مشخصهٔ دوربین
در ساده‌ترین سناریوی ممکن انجام شده است،
اما در عین حال دامنهٔ کاربرد نتایج
را به این سناریو محدود می‌کند.

سه محدودیت دیگر نیز که از مطالعهٔ کنترل‌شدهٔ نهایی برمی‌خیزند
باید به‌صراحت ذکر شوند.
نخست، تنوع پارامترهای دوربین در آزمایش‌های نهایی
از طریق افزونه‌سازی هندسی (برش‌های مقیاس‌متغیر با $K$ اصلاح‌شده)
شبیه‌سازی شده است، نه با داده‌های واقعی چنددوربینه؛
اگرچه این شبیه‌سازی از منظر هندسی دقیق است،
اثرات ثانویهٔ دوربین‌های واقعی
(نظیر اعوجاج لنز، عمق میدان و نویز حسگر)
را پوشش نمی‌دهد.
دوم، در دو دور متریکِ مبتنی بر زوم‌ـ‌برش، برتری \proposedmodel{}
در RMSE تثبیت شد، اما AbsRel دور دوم در آستانهٔ معناداری بود
(ویلکاکسون $p=0.040$ و آزمون t زوجی $p=0.089$).
این محدودیت مربوط به آن دو دور است؛ در آزمایش لرزش کانونی
تمام معیارها با $p\leq10^{-24}$ معنادار شدند. بنابراین قدرت شواهد
به میزان ابهام هندسی ایجادشده در داده وابسته است.
سوم، مقایسه‌های کنترل‌شدهٔ نهایی بر یک محک درون‌دامنه‌ای
(\lr{NYUv2}) متمرکز بوده‌اند؛
تعمیم سهم پارامترهای دوربین به فاین‌تیون بر محک‌های دیگر
(به‌ویژه صحنه‌های بیرونی با دامنهٔ عمق بزرگ)
نیازمند آزمایش‌های تکمیلی است.
در صورت دسترسی به دادهٔ آموزشی واقعی \lr{KITTI}، تکرار مطالعهٔ کنترل‌شدهٔ متریک بر روی آن می‌تواند تعمیم‌پذیری یافته‌ها را در محیط بیرونی و دامنهٔ عمق بزرگ مستند کند.

چهارم، هر بازوی هر زوج آموزشی تنها یک بار آموزش داده شده است. قفل بذر، ترتیب داده و افزونه‌سازی را میان دو بازو یکسان می‌کند و آزمون‌های زوجی عدم قطعیت در سطح تصویر را می‌سنجند، اما جایگزین تکرار کامل آموزش با چند بذر مستقل نیست. ازاین‌رو، استنباط علّی درون هر زوج قوی است، ولی برآورد واریانس ناشی از تصادفی‌بودن کل فرآیند آموزش همچنان به آزمایش تکمیلی نیاز دارد.

با در نظر گرفتن این موارد،
می‌توان گفت که محدودیت‌های مطرح‌شده
نتایج اصلی پژوهش را نقض نمی‌کنند،
بلکه چارچوب اعتبار آن‌ها را مشخص می‌سازند
و زمینه‌ای مناسب
برای توسعه و تکمیل پژوهش
در مطالعات آینده فراهم می‌کنند.


جدول~\ref{tab:limitations_future} این محدودیت‌ها را در کنار اثر هر یک بر
دامنهٔ اعتبار یافته‌ها و مسیر پژوهشی متناظر برای رفع آن قرار می‌دهد.

\begin{table}[htbp]
\centering
\caption[نگاشت محدودیت‌های پژوهش به مسیرهای پژوهشی آینده]{نگاشت محدودیت‌های
شناسایی‌شدهٔ پژوهش به اثر هر یک بر دامنهٔ اعتبار یافته‌ها و مسیر پژوهشی
متناظر برای رفع آن.}
\label{tab:limitations_future}
\small
\renewcommand{\arraystretch}{1.3}
\begin{tabularx}{\textwidth}{@{}
  >{\raggedleft\arraybackslash}p{0.20\textwidth}
  >{\raggedleft\arraybackslash}X
  >{\raggedleft\arraybackslash}X@{}}
\toprule
محدودیت & اثر بر دامنهٔ اعتبار یافته‌ها & مسیر پژوهشی پیشنهادی \\ \midrule
وابستگی به صحت ماتریس $K$ &
همان اتکای علّی که مزیت مدل است، آن را به خطای کالیبراسیون حساس می‌کند &
مطالعهٔ رفتار مدل تحت عدم قطعیت‌های واقع‌گرایانهٔ کالیبراسیون؛ طراحی سازوکار مقاوم به نویز $K$ \\
شبیه‌سازی تنوع دوربین با افزونه‌سازی &
اثرات ثانویهٔ دوربین‌های واقعی مانند اعوجاج لنز، عمق میدان و نویز حسگر پوشش داده نشده است &
آموزش بر داده‌های \emph{واقعاً} چنددوربینه؛ الگوی پاسخ به دوز پیش‌بینی می‌کند سهم $K$ در این حالت بزرگ‌تر شود \\
تمرکز بر یک محک درون‌دامنه‌ای &
تعمیم یافته‌ها به صحنه‌های بیرونی با دامنهٔ عمق بزرگ آزموده نشده است &
تکرار مطالعهٔ کنترل‌شدهٔ متریک روی \lr{KITTI} یا محک‌های بیرونی مشابه \\
تک‌اجرا بودن هر بازو &
استنباط علّی درون هر زوج قوی است، اما واریانس ناشی از تصادفی‌بودن کل فرآیند آموزش برآورد نشده است &
تکرار هر زوج با چند بذر مستقل و گزارش میانگین و انحراف معیار میان اجراها \\
تک‌معماری بودن ارزیابی &
نتایج در چارچوب خانوادهٔ \lr{Depth Anything V2} تفسیرپذیرند &
پیاده‌سازی همان چارچوب روی انکودرها و دیکودرهای متفاوت \\
عدم پیمایش محل تزریق توکن &
تزریق تنها در نخستین لایهٔ انکودر آزموده شده است &
پیمایش نظام‌مند عمق تزریق و بررسی تزریق چندلایه \\
نبود پروتکل کیفی قفل‌شده &
نتیجه‌گیری به شواهد کمی و آزمون‌های علّی محدود مانده است &
قفل‌کردن شناسهٔ تصاویر پیش از مشاهده و ذخیرهٔ خروجی هر دو بازو با نگاشت رنگ و بازهٔ عمق یکسان \\
\bottomrule
\end{tabularx}
\end{table}

% sec 6
\section{پیشنهادهایی برای تحقیقات آینده}
\label{sec:chapter5_future_work}

یافته‌ها و محدودیت‌های مطرح‌شده در این پژوهش،
زمینه‌های روشنی را برای ادامه و توسعهٔ تحقیقات آینده
در حوزهٔ تخمین عمق تک‌دوربینه فراهم می‌کنند.
پیشنهادهای ارائه‌شده در این بخش
به‌طور مستقیم بر چارچوب مفهومی و عملی
تحقیق حاضر استوار هستند
و از طرح موضوعات کلی و خارج از دامنهٔ پژوهش
اجتناب شده است.

یکی از مسیرهای طبیعی برای توسعهٔ این پژوهش،
بررسی رفتار \proposedmodel{}
در شرایطی است که ماتریس مشخصهٔ دوربین
به‌صورت دقیق در دسترس نیست.
در بسیاری از کاربردهای واقعی،
پارامترهای درونی دوربین
ممکن است همراه با نویز
یا به‌صورت تخمینی ارائه شوند.
مطالعهٔ پایداری \proposedmodel{}
در برابر خطاهای موجود در ماتریس $K$
و یا طراحی سازوکارهایی
برای مقابله با این عدم قطعیت،
می‌تواند گام مهمی در جهت کاربردپذیری عملی‌تر
روش ارائه‌شده باشد.

مسیر پژوهشی دیگر،
یادگیری هم‌زمان عمق و پارامترهای دوربین است.
در چارچوب تحقیق حاضر،
ماتریس مشخصهٔ دوربین
به‌عنوان ورودی معلوم در نظر گرفته شده است.
بررسی امکان تخمین یا اصلاح پارامترهای درونی دوربین
به‌صورت هم‌زمان با تخمین عمق،
می‌تواند منجر به مدل‌هایی شود
که وابستگی کمتری به کالیبراسیون اولیه دارند
و در سناریوهای متنوع‌تری قابل استفاده هستند.
این مسیر پیشینهٔ روشنی دارد:
\lr{UniDepth} یک شاخهٔ دوربین را در کنار شاخهٔ عمق آموزش
می‌دهد~\cite{Piccinelli2024UniDepth} و \lr{Depth Pro} فاصلهٔ کانونی را
مستقیماً از تصویر تخمین می‌زند~\cite{Bochkovskii2025DepthPro}؛
مزیت طرح کنترل‌شدهٔ این پژوهش آن است که می‌توان سهم $K$ تخمینی را در برابر
$K$ مرجع، با همان بازوی کنترل هم‌داده سنجید.

مسیر مهم دیگری که مستقیماً از یافته‌های مطالعهٔ کنترل‌شده برمی‌خیزد،
آموزش بر روی داده‌های \emph{واقعاً} چنددوربینه است.
الگوی پاسخ به دوز مشاهده‌شده نشان می‌دهد که
سهم پارامترهای دوربین با میزان تنوع هندسی داده رشد می‌کند؛
بنابراین انتظار می‌رود آموزش بر ترکیبی از مجموعه‌داده‌هایی
با دوربین‌ها و میدان‌های دید متفاوت
(به‌جای شبیه‌سازی تنوع از طریق افزونه‌سازی)
این سهم را بزرگ‌تر و پایدارتر آشکار کند؛
همان رژیمی که در آن، تناقض برچسب‌های متریک میان دوربین‌های مختلف
انگیزهٔ اصلی طراحی فضای دوربین متعارف در \lr{Metric3D} بوده
است~\cite{Yin2023Metric3D}.
در همین راستا، مطالعهٔ نظام‌مند محل تزریق توکن دوربین
در عمق‌های مختلف انکودر
(از طریق پارامتر \lr{\texttt{-{}-cam-\allowbreak token-\allowbreak inject-\allowbreak layer}})
و بررسی سازوکارهای تزریق چندلایه،
می‌تواند اثرگذاری اطلاعات هندسی را افزایش دهد.

از منظر استحکام آماری، تکرار هر زوج \proposedmodel{}/\controlmodel{} با چند بذر مستقل و گزارش میانگین و انحراف معیار میان اجراها ضروری است. این تکرار باید در کنار آزمون‌های زوجی سطح تصویر انجام شود تا دو منبع عدم قطعیت --- تفاوت دشواری تصاویر و تصادفی‌بودن آموزش --- به‌صورت جداگانه اندازه‌گیری شوند.

همچنین،
تعمیم چارچوب پیشنهادی
به سناریوهای چندتصویری یا ویدئویی
می‌تواند زمینهٔ پژوهشی ارزشمندی باشد.
اطلاعات زمانی یا چندنمایی،
قیود هندسی مکملی نسبت به اطلاعات تک‌تصویری فراهم می‌کنند
و ترکیب آن‌ها با ادغام صریح ماتریس مشخصهٔ دوربین
می‌تواند به تخمین‌های پایدارتر و دقیق‌تر
از عمق متریک منجر شود.

از منظر معماری،
بررسی پیاده‌سازی روش پیشنهادی
بر روی انکودرها و دیکودرهای متفاوت
و تحلیل میزان وابستگی نتایج
به انتخاب معماری پایه
می‌تواند به تعمیم‌پذیری بهتر چارچوب ارائه‌شده کمک کند.
چنین بررسی‌هایی امکان ارزیابی این موضوع را فراهم می‌کنند
که آیا مزایای مشاهده‌شده
یک ویژگی عمومی رویکرد هندسی پیشنهادی هستند
یا به معماری خاصی وابسته‌اند.
این مسیر، آزمون مستقیم فرضیه‌ای است که در
بخش~\ref{subsec:protocol_capacity_dependence} مطرح شد:
سهم $K$ با افزایش ظرفیت انکودر کاهش می‌یابد.
تا آنجا که بررسی شد، هیچ‌یک از کارهای شرطی‌شده بر هندسهٔ
دوربین~\cite{Facil2019CAMConvs,Yin2023Metric3D,Piccinelli2024UniDepth}
این وابستگی را به‌صورت جداگانه گزارش نکرده‌اند.

در مجموع،
پیشنهادهای مطرح‌شده در این بخش
نشان می‌دهند که
پژوهش حاضر می‌تواند
به‌عنوان نقطهٔ آغاز
برای مجموعه‌ای از مطالعات عمیق‌تر
در زمینهٔ ترکیب هندسهٔ تصویربرداری
و یادگیری عمیق در تخمین عمق تک‌دوربینه
مورد استفاده قرار گیرد.

\section{جمع‌بندی نهایی فصل}
\label{sec:chapter5_final_summary}

در این فصل، جمع‌بندی نهایی پژوهش حاضر
با تمرکز بر برداشت‌های حاصل از یافته‌ها ارائه شد.
بر اساس نتایج نظری و تجربی به‌دست‌آمده،
نشان داده شد که
تخمین عمق متریک پایدار در چارچوب تک‌دوربینه
بدون لحاظ صریح هندسهٔ دوربین
به‌صورت ذاتی با محدودیت مواجه است.
این برداشت،
ریشه در اصول هندسهٔ تصویربرداری دارد
و با نتایج عملی گزارش‌شده در فصل‌های پیشین
هم‌راستا است.

یافته‌های این پژوهش نشان می‌دهد که
ادغام آگاهانهٔ ماتریس مشخصهٔ دوربین
در فرآیند تخمین عمق
می‌تواند نقش تعیین‌کننده‌ای
در کاهش ابهام مقیاس
و افزایش پایداری متریک خروجی ایفا کند.
این بهبودها نه‌تنها در معیارهای ارزیابی،
بلکه در کیفیت بصری نقشه‌های عمق
و رفتار مدل در سناریوهای متنوع تصویربرداری
قابل مشاهده بودند.

در مجموع،
پژوهش حاضر نشان می‌دهد که
ترکیب یادگیری عمیق با قیود هندسی کلاسیک،
مسیر مؤثری برای ارتقای تخمین عمق تک‌دوربینه است.
نتایج ارائه‌شده
می‌توانند به‌عنوان مبنایی
برای توسعهٔ مدل‌های آتی
و بازنگری در شیوهٔ ارزیابی
روش‌های تخمین عمق مورد استفاده قرار گیرند.
بدین ترتیب،
این پایان‌نامه با ارائهٔ یک چارچوب
نظری و عملی منسجم،
پاسخی روشن به مسئلهٔ مطرح‌شده در ابتدای تحقیق ارائه می‌دهد
و زمینه را برای پژوهش‌های تکمیلی
در این حوزه فراهم می‌سازد.
